\documentclass[12pt]{article}
\usepackage[utf8]{inputenc}
\usepackage[T1]{fontenc}
\usepackage{hyperref}
\title{Southwestern Bontok Language Resources for K-12 Curriculum Development}
\author{Compiled by Claude Sonnet 5 (claude-sonnet-5) for Andr\'as Kornai}
\date{}
\begin{document}
\maketitle

\section{Introduction}
Southwestern Bontok (Southwestern Bontoc, ISO 639-3 vbk) is an Austronesian language of Bontoc Municipality, Mountain Province, Philippines (around Alab, Balili, Gonogon and nearby villages), and part of the Bontok macrolanguage. Search reported about 2,470 speakers in 2007, mostly older people, with a shift towards Central Bontok, Ilocano, Filipino and English. Reported dialects are Ina-ab, Binalili and Ginonogon. Search found no educational resources specific to this variety.

\section{Primary Educational Resources}
None specific to Southwestern Bontok were found in search results.

\section{Linguistic and Reference Materials}
\begin{itemize}
\item Wikipedia, ``Southwestern Bontoc''. \url{https://en.wikipedia.org/wiki/Southwestern_Bontoc}
\item Wikipedia, ``Bontoc language'': the macrolanguage and its varieties. \url{https://en.wikipedia.org/wiki/Bontoc_language}
\end{itemize}
Materials for other Bontok varieties may exist but were not examined and should not be assumed to apply to this one.

\section{Religious and Community Resources}
Joshua Project reports that Scripture is available to this group only in related Bontoc varieties and in Filipino, not in this variety. \url{https://joshuaproject.net/people_groups/20531}

\section{Digital and Audio Resources}
Global Recordings Network lists this language but reports no recordings available. \url{https://globalrecordings.net/en/language/vbk}

\section{Curriculum Development Context}
The language is endangered, spoken mainly by older adults, and is not known to be taught in schools. Any curriculum would have to begin with documentation and community-led material, and could draw on wider Bontok-language work only after checking how far it fits this variety.

\end{document}